\chapter{Error Prevention}\label{sec:ErrorPrevention}
While the whole book is more or less dealing with strategies to avoid errors in your code, this chapter is focussing on some techniques that should help to prevent error conditions to appear after your program was shipped.
 
In chapter \tqfullvref{sec:ErrorHandling} we discussed how to deal with errors (failures) when they occur. Some errors are caused by the environment (missing files, dead networks, unavailable external systems), others are caused by inappropriate handling of the software (illegal input values, wrong configurations, missing credentials). These problems are mostly unavoidable, at least from the view as a programmer. Of course you can design the UI of your product that it guides the users as much as possible to prevent them from doing things wrong, but at the very end, they are still humans (hopefully~…).

There are three other types of errors that can affect a software system:
\begin{enumerate}[nosep]
	\item{Syntax Errors}
	
	\item{Implementation or Logical Flaws}
	
	\item{Design Faults}
\end{enumerate}

For a programming language like \Java that requires a compilation step before a program can be executed, errors can be generally sorted into two categories: into compile time errors and in errors that surface during execution (usually referred to as \bdq{runtime errors}). The latter can be then put into two sub-categories: those errors that surface during testing the application, and the most nasty one, the errors that first show up in production.

Syntax errors clearly belong to the first category, and usually, they are very easy to detect and in most cases easy to fix. In fact, you \textit{must} fix them, otherwise your code will never reaches production level.

Basically, this whole book is discussing techniques that should help you avoiding errors. Those described in this chapter are explicitly targeted to this goal, while the others are more implicit. Unfortunately is it difficult to provide recipes for design errors, as they often surface only after a long time (and often the design was originally right at the time it was made, but other design decisions made it wrong later).

Consequently, we are focussing to implementation and logical errors for the remainder of this chapter.

%------------------------------------------------------------------------------

\section{Swap Logic Errors for Compiler Errors}\label{sec:SwapToCompilerErrors}
The usual weapon to fight implementation and logical errors are (unit\index{unit test}) tests\index{test}. How to write good tests is an art in its own and deserves another book, but some basics are discussed in chapter \tqfullvref{sec:WritingTests}.

But some of these errors can be prevented by swapping them to compiler errors. From a book about Java programming\index{Simmons, Robert Jr.}
\begin{quote}
“One fundamental principle of programming is that, generally, it is best to swap a logic error for a compiler error. Compiler errors tend to be found in seconds and are corrected just as fast. Syntax errors are a good example […]

Logic errors, on the other hand, are the bane of all programmers. They hide and hate to reveal themselves. Logic errors seem to have minds of their own, constantly evading detection and dodging your efforts to pin down their cause. They can easily take a thousand times more effort to solve than the worst compiler errors. Worst of all, many logic errors are not found at all and occur only intermittently in sensitive places, which causes your customer to scream for a fix. Logic errors often require you to throw thousands of man-hours at them, only to finally discover that they are minor typos.”

Robert Simmons Jr.: \textit{Hardcore Java} \autocite{Simmons:HARDCORE_JAVA}
\end{quote}

A quite simple example for swapping a potential logic error for a compiler error is the recommendation to write comparisons always with the unchangeable value on the left side. Originally, this recommendation originates from C/C++ programming where \lstinline|if| conditions are checking values of type integer – and in C, nearly everything can be an integer, or at least be interpreted like one. Java forces that the type of the expression in the \lstinline|if| condition has to be \lstinline|boolean|, therefore this approach is less useful for Java code. 

When you follow this recommendation, and you forgot the second equal sign for a comparison on equal, the compiler will complain immediately\footnote{Modern IDEs\index{IDE} will also issue warnings about assignment in \lstinline|if| clauses, even when they are syntactically correct.}
\begin{lstlisting}[language=C]
if( length() == len ) …
if( 5 == len ) …
\end{lstlisting}

Another sample for this is the recommendation to name fields with the \bdq{\lstinline|m\_|} prefix (refer to chapter \tqfullvref{sec:AttributeNames}) instead of using the \bdq{\lstinline|this.|} prefix. The compiler will never complain if you omit that prefix when accessing a field, and it will not complain if you name a local variable in the same way as a field – but it will scream loudly about a non-existing reference if you forget the \bdq{\lstinline|m\_|} prefix when accessing a field named along the rules defined here. Of course you can configure both the Eclipse IDE\index{Eclipse!IDE} and IntelliJ IDEA\index{IntelliJ IDEA} to raise an error or a warning if you access a field without \bdq{\lstinline|this.|}, or if you shadow a field through a local variable with the same name, but this is not mandatory, and therefore not completely reliable. Nevertheless, I prefer the \bdq{\lstinline|m\_|} prefix for fields because it is also easier to read.

\chapter*{Road Work ahead!}

\subsubsection{Principles}
\begin{enumerate}[label=P\arabic*.]
\item{Utilise the capabilities of your development tools to detect bugs as soon as possible.

Java is a strongly typed language, and this makes the compiler a powerful tool in the regard, given that you design your code in a manner that it can detect errors.}

\item{Take warnings serious!

The \verb#javac# in its standard configuration emits already several warnings, build tools like Maven\autocite{APACHE_MAVEN}\index{Apache Maven}\index{Maven|see{Apache Maven}} or Gradle\autocite{GRADLE}\index{Gradle} can add to them, and finally you can configure additional warnings in your favourite IDE.}

\item{Avoid to outsmart the compiler!

Casting an object to another type is a sample for this, invoking methods through Reflection\index{reflection} is another.\footnote{Thanks to the creators of Java that you cannot invoke a method through a pointer, as it is non uncommon in C and C++~…}}
\end{enumerate}


Also defining and using value types and domain value types instead of simple types can prevent logic errors from remaining undetected: providing an article number to a method that expects a customer id will not work, when you defined the domain value types \lstinline|ArticleId| and \lstinline|CustomerId| and declared the respective method to take an argument of type \lstinline|CustomerId| instead of a simple \lstinline|String| or \lstinline|int|. Refer to the chapter \tqfullvref{sec:ValueTypes} for more details about this topic.

For the same reason you should always consider to use enums instead of defining magic numbers\index{magic number}.

The best thing about compiler errors is, that they usually show up early during development – and not first after deployment at customer side. Under this aspect you should take compiler warnings seriously, no matter whether emitted by \verb#javac# or your development environment. When CI/CD pipelines are used, additional code checking tools may have been integrated. Utilise them, and respond to their warnings, at least by thinking twice before adding the \lstinline|@SuppressWarnings|\autocite{ORACLE_DOC_SUPPRESSWARNINGS_ANNOTATION}\index{java.lang!SuppressWarnings} annotation.

See also chapter \tqvref{sec:CompilerWarningsAndErrors} about compiler warnings and errors.

In this regard the use of \lstinline|java.lang.Class::forName|\autocite{ORACLE_DOC_CLASS:forName}\index{java.lang!Class!forName()} and the use of Reflection\index{reflection} in general are another nasty source for bugs that might be difficult to identify.

Basically, \lstinline|Class.forName()|\index{java.lang!Class!forName()} takes a string as its argument that is the fully qualified name of a class. If there is at least one instance of \lstinline|Class|\index{java.lang!Class} on the current classpath, the method will return that object – otherwise it will throw a \lstinline|ClassNotFoundException|\autocite{ORACLE_DOC_CLASSNOTFOUNDEXCEPTION_CLASS}\index{java.lang!ClassNotFoundException}. This will not happen at compile time, but only during the execution of the program, depending on the execution path of your program – meaning the bug that you introduced by having a typo in the class or package name may remain undetected for a long time.

So if you want to access a class object that you expect to be there, always use “\lstinline|<ClassName>.class|”\index{java.lang!Class!class} where \verb#<ClassName># is the name of the class; you can use the fully qualified name, or you can import it. Either way, the compiler will complain if the class is not there.

Only if a class is optional for your application (for example, it may come with a plugin), use “\lstinline|Class.forName( <ClassName> )|”\index{java.lang!Class!forName()}; of course you also have to add the proper error handling for the case that the class is not there.

For Reflection\index{reflection}, the situation is similar: you refer to the class element that you want to get access to by its name, provided as a string. The compiler is unable to validate that name, so a potential error will surface earliest at runtime. Refer to the chapter \tqfullvref{sec:Reflection} for more details about Reflection\index{reflection} and its alternatives.

%==============================================================================
% End of file!!
%==============================================================================



%==============================================================================
% End of file!!
%==============================================================================

\chapter*{Roadwork ahead!}

\keeplisting{10}
Another example from this category is this one:
\index{java.lang!String}\index{java.nio.charset!UnsupportedEncodingException}
\begin{lstlisting}
final String s;
try
{
    s = new String( byteArray, "UTF-8" );
}
catch( final UnsupportedEncodingException e )
{
    throw new ImpossibleExceptionError( "UTF-8 must exist!", e );
}
\end{lstlisting}

The documentation for \lstinline|java.nio.charset.Charset|\index{java.nio.charset!Charset} contains a list of standard charsets \autocite{ORACLE_DOC_CHARSET:StandardCharsets} that must be supported by each JDK/JVM implementation. According to that list, the UTF-8\index{UTF-8} \autocite{WIKIPEDIA:UTF8} encoding is mandatory and will be availabe always, and consequently, the code sequence above will \textit{never} throw an \lstinline|UnsupportedEncodingException|. 

Nevertheless, exceptions\index{exception} should never be swallowed. Instead we throw an error in the \lstinline|catch|\index{throw!catch} block, indicating the error condition. In the improbable case that UTF-8\index{UTF-8} is no longer supported by an upcoming \Java runtime environment, the problem will surface, and the program will not die silently without any hint about the reason. It also protects you from malicious changes to the name of the character encoding: a none existing \lstinline|Charset|\index{java.nio.charset!Charset} will cause noise.

Of course, here the issue could be addressed differently, by using an instance of \lstinline|java.nio.charset.Charset|\index{java.nio.charset!Charset} \autocite{ORACLE_DOC_CHARSET_CLASS} as the argument to the \lstinline|String|\index{java.lang!String} constructor, instead of the \textit{name} of the character encoding. For this case, the constructor does not declare an \lstinline|UnsupportedEncodingException|\index{java.nio.charset!UnsupportedEncodingException}.
 
\keeplisting{1}
That would look like this:
\index{java.lang!String}\index{java.nio.charset!StandardCharsets}
\begin{lstlisting}
final String s = new String( byteArray, StandardCharsets.UTF_8 );
\end{lstlisting}

Now already the compiler will complain when UTF-8\index{UTF-8} is no longer supported or the name was changed.\footnote{See also chapter \tqfullvref{sec:ErrorPrevention}.}
